\chapter{属性系统的动态变换：场地、结界、加强与反转}\label{chap:dynamic-attr}

本章只保留当前已经和源码主链直接对应、且对战斗结果有明确影响的四类机制：战场级的属性场地、单位级的四属性结界、职业技能中的属性加强，以及用标志位控制的属性反转。它们进入战斗公式的位置并不相同，因此不能混成同一类“属性状态”来理解。

\section{属性系统的整体框架}

若把一个单位在时刻 $t$ 的属性相关状态记为
\begin{equation}
    Z_t = \bigl(F_t,\;B_t,\;R_t,\;A_t\bigr),
\end{equation}
其中 $F_t$ 表示场地状态，$B_t$ 表示四属性结界，$R_t$ 表示属性反转标志，$A_t$ 表示当前生效的四属性值与抗性，那么战斗回合中的属性系统可以看成一个离散时间状态转移系统：
\begin{equation}
    Z_{t+1} = \mathcal{U}(Z_t, C_t, H_t),
\end{equation}
其中 $C_t$ 表示本回合动作集合，$H_t$ 表示本回合命中结果集合，$\mathcal{U}$ 则表示回合结算、状态写入、回合衰减与重算属性的联合算子。

这个统一写法的价值在于：它提醒我们“属性系统”不是一条公式，而是一组在不同阶段读写的状态。场地在伤害层被读取，结界在防守端减伤层被读取，属性加强若存在则应在属性重算或伤害前层进入，属性反转则在预处理重算属性时被读取。因此，若不先区分状态读写阶段，就很难解释为什么一些状态表现为覆盖，另一些却表现为开关翻转。

\section{属性场地与场地倍率}

当前真正的属性场地机制并不是职业共通技能“自我强化”，而是战场结构 \code{BattleArray[battleindex]} 上保存的场地属性、场地强度和持续回合数。若场地属性记为 $a_f$，场地强度记为 $p_f$，持续回合记为 $T_f$，则源码把它们分别写入
\begin{equation}
    \code{field_att}=a_f,\qquad \code{att_pow}=p_f,\qquad \code{att_count}=T_f.
\end{equation}

对任一单位，若其对应属性值记为 $u$，则场地倍率 \code{FieldPow} 可整理为
\begin{equation}
    \Phi(u,p_f)=0.5+\frac{up_f}{20000}.
    \label{eq:fieldpow}
\end{equation}
由于单位属性值与场地强度都不超过 $100$，所以
\begin{equation}
    0.5 \le \Phi(u,p_f) \le 1.0.
\end{equation}
当前属性伤害在经过属性相克计算之后，还会进一步乘上攻守双方场地倍率之比：
\begin{equation}
    damage_{\mathrm{field}}
    = damage_{\mathrm{attr}}
      \frac{\Phi(u_a,p_f)}{\Phi(u_d,p_f)}.
    \label{eq:field-damage}
\end{equation}
因此，场地不是单纯“给攻击方加伤”，而是通过攻守双方相对匹配度共同决定最终增减幅。

场地回合更新相对简单：每回合把 \code{att_count} 减一，当其减到非正时清空场地属性。也就是说，场地是一个战场级、单实例、按回合整体衰减的状态。

\section{四属性结界机制}

四属性结界是单位级状态，但它的存储方式与普通布尔状态不同。源码用
\begin{equation}
    \code{MAKE2VALUE(power,turn)}
\end{equation}
写入四个工作变量：
\code{CHAR_WORKFIXEARTHAT_BOUNDARY}、
\code{CHAR_WORKFIXWATERAT_BOUNDARY}、
\code{CHAR_WORKFIXFIREAT_BOUNDARY}、
\code{CHAR_WORKFIXWINDAT_BOUNDARY}。
其中高 16 位存强度 $p_b$，低 16 位存剩余回合 $T_b$。

当前实现中，结界持续回合与强度都由技能等级分档决定。若只按源码字面逻辑整理，则持续回合近似为
\begin{equation}
    T_b=
    \begin{cases}
        5, & skill\_level \ge 10,\\
        4, & skill\_level > 9,\\
        3, & skill\_level > 8,\\
        2, & skill\_level > 4,\\
        1, & \text{否则},
    \end{cases}
\end{equation}
这一分段本身存在重叠，因此更可靠的理解方式是“以代码实际执行顺序为准”，而不是按技能描述文本反推。强度则写成
\begin{equation}
    p_b=
    \begin{cases}
        100, & skill\_level \ge 100,\\
        90,  & skill\_level > 95,\\
        80,  & skill\_level > 90,\\
        70,  & skill\_level > 85,\\
        60,  & skill\_level > 80,\\
        50,  & skill\_level > 60,\\
        40,  & skill\_level > 40,\\
        30,  & skill\_level > 20,\\
        20,  & \text{否则}.
    \end{cases}
    \label{eq:boundary-power}
\end{equation}

更关键的是结界写入逻辑。若施放的是普通结界而不是“破结界”，则系统会先清掉目标侧全部四种结界，再把所选属性结界写进去。也就是说，四属性结界在当前实现里不是并存叠加关系，而是\emph{单属性覆盖关系}。若施放的是“破结界”，则会先按技能等级给一个破界成功率，成功后直接把该侧四种结界全部清空。

结界减伤公式本身也非常值得注意。若守方存在某个属性结界，则当前物理伤害分支会按攻击方对应属性值做减伤，例如土结界对应
\begin{equation}
    damage' = damage\left(1-\frac{E_a}{200}\right),
\end{equation}
其中 $E_a$ 是攻方土属性值。水、火、风结界完全类比。这里最重要的源码事实是：减伤时只检查“该结界是否存在”，实际减伤量并没有把式~\eqref{eq:boundary-power} 中存下来的 $p_b$ 再乘进去。

结界机制在实战里还有几个特别容易被误判的边界场景。第一，若同一侧已经存在某种结界，再命中另一种普通结界，则旧结界不会与新结界并存，而是先被统一清空，再由新结界覆盖写入。第二，若同一种普通结界在持续回合尚未结束时再次命中，其效果更接近“按新技能结果重写当前槽位”，而不是在旧强度、旧回合数基础上继续累加。第三，若破结界命中成功，则四种结界槽会一起清空；若失败，则原结界状态整体保留。

还可以用一个具体案例看清结界与属性反转的关系。设攻击方原始四属性记为 $(E,W,F,N)=(70,30,0,0)$，分别对应土、水、火、风；防守方身上存在土结界。若攻击方未被反转，则土结界减伤读取的是攻方当前土属性 $70$，减伤比例约为 $70/200=35\%$。若攻击方先被属性反转，则当前有效属性会变成 $(0,0,70,30)$。此时防守方虽然仍是”土结界”，但减伤读取的已经是反转后的土属性 $0$，减伤比例下降为 $0/200=0\%$。这说明结界槽位本身不会被反转交换，但它所读取的”攻方对应属性值”会随着反转后的当前属性而改变。

同理，场地倍率也应理解为读取“当回合生效属性”而不是角色的某个静态建档值。若某单位长期带有属性反转标志，则其每回合预处理后参与场地倍率计算的属性向量也会是反转后的版本。于是，属性反转不仅会改变属性相克和结界减伤，还会连带改变该单位在场地中的受益方向。

\section{属性加强机制}

这一节的结论非常简单，但必须明确写出来：当前源码中的职业共通技能“属性加强（自我强化）”对应实现仅返回 \code{TRUE}，没有看到实际数值逻辑。也就是说，就这份代码口径而言，“自我强化”的战斗内数值效果是\emph{空实现}。

这件事的重要性不在于“某个技能没写完”，而在于它提醒研究不能把玩家经验中的效果自动补回源码。若后续版本或别的分支实现了具体公式，那应当另行分析；但在当前论文口径下，属性加强只能被记为“在现有源码中未见明确数值作用”，而不能擅自假定为乘法、加法或覆盖效果。进一步说，这也意味着几个常见的边界猜想在当前代码下都不成立：多次施放属性加强不会形成层数累加，属性加强也不会与结界、场地或属性反转再叠出额外乘区，因为源码里根本没有可供叠加的数值入口。

\section{属性反转机制}

属性反转的实现是本章最需要按时序理解的部分。技能命中后，系统先对目标的 \code{CHAR_BATTLEFLG_REVERSE} 做一次异或：
\begin{equation}
    flg \leftarrow flg \oplus \code{REVERSE},
\end{equation}
然后立即调用 \code{BATTLE_AttReverse()}。后者并不是“无条件再交换一次”，而是先检查目标头上当前是否仍带有 \code{REVERSE} 标记；只有标记存在时，才把当前固定属性按如下映射交换：
\begin{equation}
    (E,W,F,N) \longmapsto (F,N,E,W),
    \label{eq:reverse-map}
\end{equation}
其中 $E,W,F,N$ 分别表示土、水、火、风。也就是说，土火互换，水风互换。

这个判定顺序会带来一个很关键的结论。第一次命中属性反转时，标记从关变开，于是 \code{BATTLE_AttReverse()} 会立刻把当前固定属性切到反转态。第二次命中同技能时，标记从开变关；这时 \code{BATTLE_AttReverse()} 进入后会直接返回，因此\emph{不会在命中当下把当前固定属性切回原态}。也就是说，“反转标记被关掉”和“当前固定属性立刻复原”在源码里不是同一件事。

真正把属性恢复成正常口径的，是后续的固定属性重算过程。在每回合预处理 \code{BATTLE_PreCommandSeq()} 中，系统会先重算本回合固定属性；只有头上仍带有 \code{REVERSE} 标记时，才会再次调用 \code{BATTLE_AttReverse()}。因此，若某单位已经处于反转态，又在同回合再次被属性反转命中，则会出现一个临时状态：\emph{标记已经被清掉，但当前这份固定属性快照仍然保持反转结果}。等到下一次固定属性重算发生时，由于标记已关，系统不再补做交换，属性才会回到正常口径。

这也正是前一版文稿里“连续两次命中就会立刻切回原态”不成立的原因。更准确的说法应当是：第二次命中会关掉反转标记，但不会立即重做一次逆交换；是否在同回合内就恢复正常，要看之后是否还有新的固定属性重算入口。若没有，就会把这份反转后的固定属性继续用到本回合后续判定里；若有，则会在那次重算后恢复成正常属性。

这一机制可以直接推出几个边际案例。第一，若单位在反转状态尚未解除时更换装备、骑宠状态改变，或本回合又被其它流程重写固定属性，则下一次重算会先按最新口径生成固定属性，再根据标记是否仍在决定要不要继续交换。因此，反转作用的是“当回合有效属性”，不是“某次命中时刻留下来的旧快照”。第二，若防守方处于反转状态而攻击方施放结界，结界写入的仍然是土、水、火、风四个固定槽位本身；被交换的是人物当前属性，而不是结界槽位的语义。换言之，反转会改变结界减伤读取到的属性值，但不会把“土结界”变成“火结界”。第三，若同一目标在同回合内先被反转、后又被同技能再次命中，就可能出现“头上已无反转标记，但当前属性仍按反转态参与余下判定”的过渡现象；这正是需要把“标记状态”和“当前固定属性快照”分开理解的地方。

从玩家语言看，常把这类效果叫作“属性转换”。本文沿用源码口径称其为“属性反转”，原因就在于它不是把某一项属性永久改写成另一项，而是通过一个可开可关的标志位，在固定属性重算之后对当前属性做映射交换。这个差异看似只是命名问题，实则决定了我们应该把它理解成一个“标记控制的重算后交换过程”，而不是一次性的数值替换。

\section{多次状态命中的属性演化}

把前几节合起来看，可以把多次状态命中分成三种完全不同的演化模式。

第一种是\emph{覆盖型}，代表就是四属性结界。普通结界命中后会先清空已有四结界，再写入新结界，因此新状态会覆盖旧状态，而不是叠加在旧状态之上。第二种是\emph{翻转型}，代表是属性反转，但这里的“翻转”严格地说只发生在标志位层；至于当前固定属性是否立刻改回去，还要看当下有没有重新进入属性重算。第三种是\emph{回合重算型}，代表属性场地、结界剩余回合以及当前固定属性本身，它们会在回合推进时统一进入新的状态快照。

从目前已核对的代码可确定：结界的“多次命中”不是强度叠加而是覆盖；反转的“多次命中”不是刷新持续时间，而是对反转标志做异或；破结界的“多次命中”则表现为若干次独立清空尝试。更细一点地说，若同回合内先后发生“普通结界命中、破结界命中、再一次普通结界命中”，则最终保留的是最后一次成功写入的普通结界；若同回合内属性反转连续命中两次，则最终的标记状态会回到关闭，但当前固定属性是否已经恢复，还要看是否发生过新的固定属性重算；若同回合内既有反转又有结界，则结界写入的是槽位，反转影响的是之后读取这些槽位时所对应的当前属性值。这些都是典型的“顺序敏感”案例。

\section{动态属性系统的状态转移模型}

综合前述各节，当前属性系统最适合写成由若干算子组成的离散更新链。若把场地算子记为 $\mathcal{F}$，结界写入算子记为 $\mathcal{B}$，反转标志与属性交换的联合算子记为 $\mathcal{R}$，回合衰减算子记为 $\mathcal{D}$，则一回合内的属性状态更新可抽象为
\begin{equation}
    Z_{t+1}
    = \mathcal{D}\circ \mathcal{R}^{\mathbf{1}_{\mathrm{reverse}}}
      \circ \mathcal{B}^{\mathbf{1}_{\mathrm{boundary}}}
      \circ \mathcal{F}^{\mathbf{1}_{\mathrm{field}}}(Z_t),
\end{equation}
其中各指标函数表示本回合是否有相应状态命中。虽然这只是抽象写法，但它足以揭示三个关键结论。

第一，属性系统不是交换群结构，不同状态写入顺序不可随意交换。例如结界是在防守端减伤层生效，而反转是在固定属性重算后生效，因此“先反转再上结界”和“先上结界再反转”不一定等价。第二，反转的可逆性成立于“标志位层”和“完整重算后的新属性快照层”，不能粗暴理解成“连续命中两次就立刻把当前值切回去”。第三，所谓“属性玩法”并不是单一伤害乘法，而是一套由战场级状态、单位级状态和标志位共同构成的动态系统。这一视角也正是后文把宠物培养、骑宠速度、法术命中与伤害统一到同一研究框架中的关键桥梁。
